\chapter*{Conclusion}
\addcontentsline{toc}{chapter}{Conclusion}

This thesis developed and applied a reproducible pipeline for population-level
characterisation of post-infarction scar geometry using standardised polar maps. Applied to
105 patients with usable core-zone segmentations from the DEVELOP-VT dataset, the
pipeline enabled scar location, segment and ring burden, transmurality, and geometric
descriptors to be compared in a common anatomical frame. The analysis showed marked
inter-patient heterogeneity, non-uniform scar distribution across ventricular regions, and the
feasibility of retrieval-based generation as a conservative, physiologically grounded baseline.

\section*{Research Questions and Findings}
\addcontentsline{toc}{section}{Research Questions and Findings}

Three research questions guided this work, each representing an attempt at a progressively
more conservative form of geometric representation.

\paragraph{RQ1: Is there a low-dimensional space that can reproduce plausible scar geometries?}

The original aim was to learn a compact latent representation of scar geometry---a
low-dimensional manifold from which new, plausible scar configurations could be sampled
and modified. This proved infeasible. The scar-pattern space is intrinsically high-dimensional:
six PCA components are required to explain only 72\% of segment-level variance at the
coarse resolution of AHA-17 segments, with PC1 accounting for only 18\%. This low
concentration of variance in the leading component is not a property of the chosen
representation---it reflects the genuine independence of infarct location, extent,
transmurality, and longitudinal distribution across patients, which vary as a function of which
vessel occludes, at what level, and for how long. A cohort of 105 patients is fundamentally
underpowered to learn a reliable latent space over a distribution this high-dimensional. In
this cohort and representation, the data did not support a sufficiently compact
low-dimensional generative space. This likely reflects the genuine biological and anatomical
heterogeneity of post-infarction scar rather than a limitation of the chosen method.

\paragraph{RQ2: Can scar geometry be described by a known parametric distribution?}

If a full generative latent space is out of reach, a weaker goal is to describe scar geometry
through a known parametric family from which new cases could be sampled. This also
proved infeasible. At every level of aggregation---individual AHA segments, coronary
territories, global burden---the distributions are zero-inflated, heavily right-skewed, and
non-normal, with mean-to-median ratios of 2.4--3.7 and standard-deviation-to-mean ratios of
2--4 at the segment level. The marginal distributions of individual descriptors can be
characterised non-parametrically, but the joint distribution across descriptors cannot be
reproduced from them: territory assignment, scar location, longitudinal gradient,
transmurality, and fragment shape are not independent, and the dependencies between them
are complex and patient-specific. Encoding a patient as a feature vector captures marginal
properties but loses the geometric relationships between them. Two patients with identical
feature vectors can present fundamentally different scar geometries. In this cohort, the data did not support a tractable (computationally and statistically manageable)
parametric description of the joint scar geometry distribution. This sharpens the problem: what is needed is not a generative
model but a representation that is honest about what it encodes and what it discards.

\paragraph{RQ3: How can post-infarction scar geometry be parameterised in a compact representation that preserves clinically meaningful structure?}

This is the question the thesis answers. The polar-map representation projects the
three-dimensional scar surface onto a standardised $64 \times 128$ bull's-eye grid anchored
to the long axis and the septal reference, providing a common $(s, \theta)$ coordinate frame
across patients regardless of ventricular size, shape, or orientation. This representation
preserves the properties that are clinically meaningful and statistically tractable: scar location
by territory and AHA segment, circumferential and longitudinal extent, transmural onset and
depth, and fragment-level shape. The population-level characterisation built on this
representation identified the mid-inferolateral and mid-inferior segments as the dominant
hotspot, confirmed the three-axis organisation of scar geometry, and showed that the
retrieval-based framework can identify geometrically similar patients with a mean segment
Dice of 0.597 under leave-one-out evaluation.

The answer to RQ3 is: polar maps, with explicit information loss. What is preserved is
enough to support population-level analysis and retrieval-based generation. What is lost is
the intrinsic three-dimensional volumetric geometry of the scar---how it develops through the
wall, how adjacent regions are connected in depth, and what its full surface morphology
looks like from the epicardial side. The large gap between segment Dice (mean 0.597) and
polar-map Dice (mean 0.106) is a partial measure of this loss at the sub-segment level; the
loss in the transmural dimension is harder to quantify but no less real.

This work approaches scar from the opposite direction to image-level synthesis methods
such as LGESynthNet \cite{jacob2024} and CLAIM \cite{ramzan2025}: it analyses scar as
a 3D geometric object rather than a pixel mask, providing the structural characterisation
layer that complements those approaches.

\section*{Reflection on Scope}
\addcontentsline{toc}{section}{Reflection on Scope}

The original ambition of this thesis was more generative in scope. The intended goal was to
learn a compact latent space of scar geometry from which new, plausible configurations could
be sampled and navigated toward arrhythmogenic substrates. This proved unreachable at the
current scale: the extreme heterogeneity of the cohort---scar varying across all coronary
territories, longitudinal levels, transmural profiles, and fragment counts---means no compact
low-dimensional structure exists that a generative model could faithfully learn from 105 cases.
The thesis therefore adopted a more analytical and reproducible strategy: characterise scar
geometry explicitly through interpretable descriptors, and use those descriptors as the basis
for retrieval-based generation. This is a more conservative approach, but it is honest about
what the data support, produces verifiable outputs---every generated scar is a real observed
case---and establishes a descriptive foundation that future generative models can build on.

\section*{Summary of Contributions}
\addcontentsline{toc}{section}{Summary of Contributions}

A cohort of 105 patients with valid left ventricular reconstructions and core-zone meshes was
assembled from 168 candidate folders spanning five acquisition years, constituting a
systematically curated and centrally organised version of the DEVELOP registry for
computational analysis.

A fully automated polar-map pipeline was implemented, with a multi-candidate
apex-detection strategy that evaluates LV PC1, RV PC1, and LV PC2 as search directions
and selects the one yielding maximum eccentricity---necessary because post-infarction
remodelling renders a single fixed direction insufficient in approximately 25\% of cases. A
continuous vertex-depth transmurality measure was developed as a replacement for the
too-coarse discrete layer meshes exported by ADAS 3D.

A population-level statistical characterisation of scar geometry was performed, establishing
non-uniform spatial distribution across AHA segments, identifying the mid-inferolateral and
mid-inferior segments (AHA 11 and 10) as the most affected regions straddling the
RCA--LCX boundary, confirming consistent basal sparing, and documenting the three-axis
geometric organisation of scar. Transmurality was characterised as predominantly
endocardial-dominant with a dense-core / thin-edge gradient.

A retrieval-based scar generation framework was proposed and evaluated under
leave-one-out cross-validation, yielding a mean absolute burden error of 0.015, mean
segment Dice of 0.597, and mean polar-map Dice of 0.106. SVD-based feature pruning
consistently degraded retrieval performance, demonstrating that discriminative information is
distributed across dimensions with low individual variance contributions.

\section*{Future Directions}
\addcontentsline{toc}{section}{Future Directions}

The most immediate extension is linkage to clinical outcome data as the DEVELOP registry
accumulates follow-up events. The feature vector developed here provides a direct interface
between population geometry and clinical prediction. On the generative side, the retrieval
framework could be conditioned on clinical variables---culprit territory, ejection fraction, time
since infarction---which the current architecture already supports by restriction of the case
base or appending of clinical features to the descriptor vector. A natural further extension
would connect the geometric representation to image-level synthesis methods such as
LGESynthNet, using the polar-map descriptors to condition image generation and produce
augmented LGE-CMR data with geometrically controlled scar properties.

The deeper open problem is the transition from the scalar-descriptor representation built
here toward a true geometric shape model of scar. One concrete pathway toward this is a
two-stage generative architecture. In the first stage, a shape encoder operating on the target
LV geometry---without scar---learns a representation of the feasible anatomical region where
scar can live, constrained by wall thickness, curvature, and coronary territory boundaries. In
the second stage, a scar geometry encoder places a scar configuration inside that feasible
region, using the polar-map descriptors and the retrieval-based library developed in this
thesis as the prior over plausible scar shapes. Decoupling these two stages is important:
trying to learn scar geometry and its anatomical context jointly is precisely what proved
infeasible at 105 cases; separating them gives each encoder a more tractable learning
problem.

The labelling problem remains: a generative model trained only on observed scars produces
plausible configurations but cannot distinguish arrhythmogenic from non-arrhythmogenic
ones. The solution is to close the loop through electrophysiological simulation. Scar
configurations generated by the two-stage architecture would be passed to a computational
model, re-entrant circuit formation would be assessed, and outcomes would be compared
against clinical endpoints. This process builds a labelled library of scar
geometries annotated by their simulated arrhythmogenic potential. With that library, a
generative model conditioned on arrhythmogenic label becomes trainable---and the latent
space becomes navigable not toward any plausible scar, but specifically toward
configurations that sustain VT. This is the pathway that connects the population geometry
characterised in this thesis to the clinical endpoint that motivated it.
